\documentclass[a4paper]{article}
\usepackage{amsfonts}
\usepackage{amsmath}
\usepackage{amssymb}
\usepackage{graphicx} 
\usepackage{cancel}

\begin{document}
  
\noindent \large Auteur: Odia Kibor \\
\noindent \large Studierichting: Informatica\\
\noindent \large Oefeningenreeks 4A nr. 42.8\\

\medskip

\normalsize


$ \displaystyle\int \dfrac{dx}{x^3-7x+6} $\\

nulpunten noemer: $1$, $2$ en $-3$\\

$ \dfrac{1}{x^3-7x+6} = \dfrac{A}{x - 1} + \dfrac{B}{x - 2} + \dfrac{C}{x + 3}$\\
 
$ A =\dfrac{1}{(x - 2)(x + 3)}\bigg|_{x = 1} = -\dfrac{1}{4} $\\

$ B =\dfrac{1}{(x - 1)(x + 3)}\bigg|_{x = 2} = \dfrac{1}{5} $\\

$ C =\dfrac{1}{(x - 1)(x - 2)}\bigg|_{x = -3} = \dfrac{1}{20} $\\

$\Rightarrow \displaystyle\int \dfrac{A}{x - 1} + \dfrac{B}{x - 2} + \dfrac{C}{x + 3} dx = \displaystyle\int -\dfrac{1}{4(x - 1)} + \dfrac{1}{5(x - 2)} + \dfrac{1}{20(x + 3)} dx$

$ = -\dfrac{\ln|x - 1|}{4} + \dfrac{\ln|x - 2|}{5} + \dfrac{\ln|x + 3|}{20} + C $\\


















\begin{thebibliography}{999}
%\bibitem{a} Referentie
\end{thebibliography}
\end{document} 


